%Plantilla para escritura de memoria/tesis
%Universidad Técnica Federico Santa María, Chile
%Departamento de Ingeniería Química y Ambiental
%Iván Cornejo www.cornejog.com/memoristas
%2021-03-01
\documentclass[12pt, oneside]{report}
\usepackage[utf8]{inputenc}
\usepackage{mathptmx}

\usepackage{titlesec}

\titleformat{\chapter}{\normalfont\Large\bf\MakeUppercase}{\thechapter.}{10pt}{\Large\bf\MakeUppercase}
\titlespacing*{\chapter}{0pt}{0pt}{40pt}

\usepackage[letterpaper,left=3.5cm,right=3cm,top=3cm,bottom=3cm]{geometry} %control de tamaño de papel y márgenes
\usepackage{graphicx} %para poder incluir figuras 
\usepackage{siunitx} %para notación científica
\usepackage{float} %complemento para figura
\usepackage{import} %para poder dividir el texto en varios archivos
\usepackage{natbib} %para la bibliografía
\usepackage{setspace} \onehalfspacing %Control del espacio entre lineas

%Hoja de la portada
\title{
\vspace{-3.0cm}
\includegraphics[width=8cm]{figuras/logo_eciem.png}\\
\vspace{3cm}
%{\Large Seminario Superior en Sistemas}\\
{\Large \textbf{TÍTULO DEL TRABAJO DE PREPARACIÓN DE SEMINARIO SUPERIOR}}\\
\vspace{1cm}
{\footnotesize \begin{flushright} \textit{Preparación de Seminario Superior parte de la actividad de} \\ \textit{Titulación para optar al Título profesional de Ingeniero} \\ \textit{en Información y Control de Gestión.} \end{flushright}}
\vspace{0.5cm}
{\footnotesize \begin{flushright} \textit{\textbf{Profesor Guía: Nombre del profesor guía}} \end{flushright}}
\vspace{1.5cm}
{\normalsize \textbf{Autor1\\Autor2}\\}
\vspace{3.0cm}
{\normalsize MES-AAA}
}
\author{} %dejar vacío
\date{} %dejar vacío
%----------

%Este bloque es para adaptar los nombres a español
\renewcommand\bibname{Bibliografía}
\renewcommand{\figurename}{Figura}
\renewcommand{\tablename}{Tabla}
\renewcommand{\listfigurename}{Índice de Figuras}
\renewcommand{\listtablename}{Índice de Tablas}
\renewcommand{\contentsname}{Índice}
\renewcommand{\chaptername}{Capítulo}
%----------

%Control del encabezado y pie de página
\usepackage{fancyhdr}
\pagestyle{fancy}
\setlength{\headheight}{15pt}
\renewcommand{\headrulewidth}{0pt}
\fancyhf{}
\rhead{}
\lhead{}
\cfoot{\thepage}
%----------

\begin{document}
\pagenumbering{roman} %comienza con número romano para las páginas
\maketitle %agrega la portada al documento

%se incluye lo contenido en resumen.tex
%\include{resumen}

%se incluye lo contenido en abstract.tex
%\include{abstract}

%se incluye lo contenido en agradecimientos.tex
%\include{agradecimientos}

%Estas 3 listas se crearán automáticamente, según los capítulos y secciones del documento 
\tableofcontents    %crea la lista de contenidos
\listoftables       %crea la lista de tablas
\listoffigures      %crea la lista de figuras

%incluye lo contenido en nomenclatura.tex
% \include{nomenclatura}

%incluye lo contenido en C1problema.tex
\include{C1problema}

%incluye lo contenido en C2marcoteorico.tex
\include{C2marcoteorico}

%incluye lo contenido en C3objetivos.tex
\include{C3objetivos}

%incluye lo contenido en C4conclusiones.tex
\include{C4medologia}

%incluye lo contenido en C5cartagantt.tex
\include{C5cartagantt}

%Generacion automatica de la lista de referencias 
\bibliographystyle{apalike} %estilo APA
%\bibliographystyle{unsrt} %estilo tipo EEE 
\bibliography{references} %indica nombre del archivo de referencias


\end{document}
